\section{Introduction}
\label{sec:introduction}

Representation engineering \citep{zou2023repe} reveals that large language models encode sycophancy as a linear direction in their residual streams.
A contrastive sycophancy direction $\dsyc$, extracted from style-controlled pairs of valid belief revision and sycophantic capitulation, achieves cross-validated AUROC $0.827 \pm 0.028$ on Qwen3-8B and $0.779 \pm 0.019$ on Llama-3-8B (\S\ref{sec:direction_quality}).
The direction captures evidential granularity: discriminability rises monotonically from weak-evidence pairs (AUROC 0.675) through medium (0.883) to strong (0.929), and the signal remains stable under LoRA training with cosine drift of approximately 0.96 at 500 steps.
These results, replicated across two architecturally distinct model families, point to an appealing possibility: if the model's representations already separate sycophantic from non-sycophantic behavior, that separation should be convertible into a reward signal for reinforcement learning.
Sycophancy is among the most persistent alignment failures in deployed systems \citep{sharma2023sycophancy}, and a representation-level reward could complement existing behavioral approaches by targeting the internal computation rather than its surface outputs.

This paper shows that the conversion fails, and that the failure runs deeper than the representation signal itself.
We test three formulations that translate $\dsyc$ into a GRPO \citep{shao2024grpo} training signal (cosine similarity, scaled cosine, and binary probe) across six experimental configurations; every formulation fails.
Cosine similarity produces a reward gap of just 0.098 between valid and sycophantic activations, too narrow to drive policy updates (\S\ref{sec:cosine_failure}).
Online recalibration improves the offline AUROC from 0.82 to 0.91 yet triggers mode collapse.
A binary probe with stopped gradients creates a stubbornness trap: the model learns that never revising its answer reliably collects the probe reward, and the true-positive flip rate ($\tof$) drops to zero within 500 steps while generation entropy remains at 0.70 (\S\ref{sec:probe_trap}).
These failures might suggest that removing $\dsyc$ in favor of a purely behavioral reward would resolve the problem.
It does not.
R3-only GRPO, which uses an NLI-based behavioral reward without any representation component, also fails to calibrate sycophancy: the \sycon{} score starts at 0.074 at checkpoint 1000, briefly matches the vanilla baseline (0.134) at checkpoint 4500, but immediately regresses to 0.114 and plateaus from step 5000 onward (\S\ref{sec:discussion}).
The ceiling of sparse behavioral reward is the untrained vanilla level---it cannot push calibration beyond the pre-training baseline.
Two layers of failure emerge: $\dsyc$-based rewards are actively harmful, and GRPO itself is structurally insufficient for sycophancy calibration.

Existing approaches to sycophancy mitigation span three categories.
Representation-based inference methods, including activation addition \citep{turner2023actadd} and contrastive activation addition \citep{rimsky2024caa}, steer model behavior at test time without modifying weights.
RL-based training methods use scalar rewards to update policy through reinforcement learning; our $\dsyc$ GRPO variants and R3-only GRPO both fall into this category, as does representation-guided RL more broadly \citep{zhang2026latentgrpo}.
SFT-based training takes a fundamentally different approach: ACT \citep{consistency2025bctact} applies supervised fine-tuning with an L2 activation consistency loss that provides dense per-token supervision rather than a single scalar per completion.
ACT achieves a \sycon{} score of 0.301 at checkpoint 8000, more than double the vanilla baseline, while every GRPO variant remains below it (\S\ref{sec:discussion}).
The contrast is informative.
Techniques for discovering latent knowledge \citep{burns2023ccs} and representation-guided RL share a premise: a direction with high discriminability in activation space should transfer into a useful training signal.
For reasoning faithfulness and safety refusal, this premise holds, because the behavior to suppress is geometrically distinct from the behavior to preserve along the relevant direction.
Sycophancy violates this premise.
Capitulation under social pressure and revision under valid evidence produce overlapping activation signatures along $\dsyc$, because both involve the model shifting its stated belief toward the user's position.

The root cause is a semantic gap between discriminability and reward utility.
$\dsyc$ encodes \emph{whether} belief changed, not \emph{whether it should have}.
A post-hoc classifier can use this signal effectively because it scores labeled examples where the validity of the belief change is already known.
A reward signal during RL training has no such access to ground truth; it must distinguish warranted revision from unwarranted capitulation using activation geometry alone, and $\dsyc$ collapses exactly this distinction.
This conflation explains both layers of failure.
$\dsyc$ rewards penalize all belief change indiscriminately, producing behavioral collapse.
R3-only GRPO's single scalar per completion cannot compensate because it lacks the per-token resolution to teach which belief changes are warranted and which are not.
ACT succeeds because its dense supervision operates at the token level, where the valid/invalid distinction can be resolved.
Sycophancy is defined relative to external evidential warrant \citep{sharma2023sycophancy}, a property that no single activation direction can fully encode.
The entanglement between helpfulness and sycophancy at the representation level, where the same mechanisms that allow models to track user intent also predispose them to capitulate \citep{wang2025internal}, is amplified by RLHF \citep{shapira2026amplifies}.
When reward overoptimization \citep{gao2023scaling} operates on a signal that conflates compliance with capitulation, the result is predictable collapse \citep{casper2023openproblems}.

We present three contributions:

\begin{itemize}
\item \textbf{All GRPO approaches fail for sycophancy calibration} (C1). We extract and validate $\dsyc$ on two model families (Qwen3-8B AUROC $0.827 \pm 0.028$; Llama-3-8B AUROC $0.779 \pm 0.019$), then show that all three $\dsyc$-based reward formulations produce behavioral collapse (\S\ref{sec:cosine_failure}, \S\ref{sec:probe_trap}). R3-only GRPO, which removes $\dsyc$ entirely, also fails: \sycon{} briefly matches vanilla (0.134) at checkpoint 4500 but regresses to 0.114 and plateaus (\S\ref{sec:discussion}). Discriminability does not confer reward utility, and sparse RL rewards are structurally insufficient for this task.

\item \textbf{Probe reward trap with behavioral-before-distributional collapse} (C2). The binary probe reward produces behavioral death ($\tof = 0$ at step 500) while entropy remains at 0.70. The 1,300-step lag before distributional collapse means standard entropy monitoring misses the failure entirely, with the reward averaging $+0.72$ throughout (\S\ref{sec:probe_trap}).

\item \textbf{Sparse RL reward cannot convey the valid/invalid distinction} (C3). We trace all failures to a unified root cause: $\dsyc$ encodes belief-change direction but not evidential warrant. A single scalar reward per completion, whether from representations or behavioral classifiers, cannot teach the model which belief changes are warranted. Dense SFT supervision resolves this: ACT achieves \sycon{} of 0.301, while no GRPO variant surpasses vanilla (\S\ref{sec:discussion}).
\end{itemize}